%% kapa: exam=1 round=36 year=2025 period=1 form=A subject=economics subject_ko=경제학원론 seq=3 number=43
%% source: Q-Net 감정평가사 기출문제 (cst00309) artlSeq=5231525 file=raw/1008/5231525/2025년 제36회 감정평가사 1차시험 1교시 문제지 원본.pdf
%% conversion: kordoc 4.24.51 -> tools/gen1.py (auto); review=proofread-sonnet; flags=none
\begin{question}{43}
코즈(R. Coase)정리가 성립하기 위한 조건으로 옳지 않은 것을 모두 고른 것은?
\begin{framedtext}
ㄱ. 정부는 외부효과의 크기를 정확하게 파악하고 있다. \par
ㄴ. 관련 당사자 간 협상이 가능하다. \par
ㄷ. 협상에 따른 거래비용이 발생하지 않는다. \par
ㄹ. 당사자 간 재산권이 명확하게 설정되어 있다.
\end{framedtext}
\begin{choices}
  \item ㄱ
  \item ㄴ
  \item ㄱ, ㄹ
  \item ㄴ, ㄷ
  \item ㄷ, ㄹ
\end{choices}
\end{question}
\answer{1}
